\Needspace{7\baselineskip}
\section{The Exact \texorpdfstring{\(206\)}
\label{sec:13}{206}-to-\texorpdfstring{\(409\)}{409} Construction}
The intrinsic middle-to-minimum ratio is greater than \(206\) and less than \(207\). This section first proves that interval without measured input. It then uses the one-event/two-preserved-face recovery record to construct the finite count \(409\). The spectrum supplies \(206\); the recovery record supplies the two-face structure needed to combine two \(206\)-element supports. Figure \(\ref{fig:observer-betweenness}\) summarizes observer-rooted middle placement, while Figure \(\ref{fig:two-face-409}\) separately summarizes free amalgamation over the common three-root carrier.

\Needspace{5\baselineskip}
\subsection{Exact Determination of \texorpdfstring{\(206\)}{206}}
With the local Taylor variable

\begin{equation*}
\begin{adjustbox}{max width=\linewidth,center}
$\displaystyle
x:=\theta_o=\frac29,
$
\end{adjustbox}
\end{equation*}
one has

\begin{equation*}
\begin{adjustbox}{max width=\linewidth,center}
$\displaystyle
\lambda_{\min}
=
1-\frac{\cos x}{\sqrt2}
-\frac{\sqrt6}{2}\sin x,
$
\end{adjustbox}
\end{equation*}
\begin{equation*}
\begin{adjustbox}{max width=\linewidth,center}
$\displaystyle
\lambda_{\mathrm{mid}}
=
1-\frac{\cos x}{\sqrt2}
+\frac{\sqrt6}{2}\sin x.
$
\end{adjustbox}
\end{equation*}
% LAYOUT_ONLY_GUARD:PELL_FIRST_CERTIFICATE_CHAIN
\Needspace{18\baselineskip}
Because \(0<x<1\), the alternating Taylor theorem gives finite rational upper and lower bounds for \(\sin x\) through degree nine and for \(\cos x\) through degree eight. The exact Pell certificates

\begin{equation*}
\begin{adjustbox}{max width=\linewidth,center}
$\displaystyle
2(470832)^2-(665857)^2=-1,
\qquad
2(1136689)^2-(1607521)^2=1
$
\end{adjustbox}
\end{equation*}
% LAYOUT_ONLY_GUARD:PELL_BOUND_GIVE
\Needspace{7\baselineskip}
give

\begin{equation*}
\begin{adjustbox}{max width=\linewidth,center}
$\displaystyle
\frac{470832}{665857}
<
\frac1{\sqrt2}
<
\frac{1136689}{1607521},
$
\end{adjustbox}
\end{equation*}
% LAYOUT_ONLY_GUARD:PELL_SECOND_CERTIFICATE_CHAIN
\Needspace{16\baselineskip}
while

\begin{equation*}
\begin{adjustbox}{max width=\linewidth,center}
$\displaystyle
4(1046629)^2-6(854569)^2=-2,
\qquad
4(4656965)^2-6(3802396)^2=4
$
\end{adjustbox}
\end{equation*}
% LAYOUT_ONLY_GUARD:PELL_BOUND_GIVE
\Needspace{7\baselineskip}
give

\begin{equation*}
\begin{adjustbox}{max width=\linewidth,center}
$\displaystyle
\frac{1046629}{854569}
<
\frac{\sqrt6}{2}
<
\frac{4656965}{3802396}.
$
\end{adjustbox}
\end{equation*}
Sign-aware rational interval propagation yields

\begin{equation*}
\begin{adjustbox}{max width=\linewidth,center}
$\displaystyle
206\lambda_{\min}^2
<
\lambda_{\mathrm{mid}}^2
<
207\lambda_{\min}^2.
$
\end{adjustbox}
\end{equation*}
Since \(\lambda_{\min}>0\),

\begin{equation}
\label{eq:206-interval}
\begin{adjustbox}{max width=\linewidth,center}
$\displaystyle
\boxed{
206
<
R_{\mathrm{mid}/\min}^{\mathrm{int}}
=
\left(
\frac{\lambda_{\mathrm{mid}}}{\lambda_{\min}}
\right)^2
<
207.
}
$
\end{adjustbox}
\end{equation}
Consequently,

\begin{equation*}
\begin{adjustbox}{max width=\linewidth,center}
$\displaystyle
\boxed{
n_{\mathrm{mid}}
:=
\left\lfloor
R_{\mathrm{mid}/\min}^{\mathrm{int}}
\right\rfloor
=206,
}
$
\end{adjustbox}
\end{equation*}
and \(\rho_{\mathrm{mid}}=R_{\mathrm{mid}/\min}^{\mathrm{int}}-206\in(0,1)\). No decimal and no measured muon value participates in this result.

\Needspace{5\baselineskip}
\subsection{Two Supports Attached to the Preserved Faces}
The carried record supplies one event, its two preserved faces, and

\begin{equation*}
\begin{adjustbox}{max width=\linewidth,center}
$\displaystyle
(n_{\mathrm{post}},n_{\mathrm{route}})=(2,1).
$
\end{adjustbox}
\end{equation*}
% LAYOUT_ONLY_GUARD:ROOT_CARRIER_SEQUENCE
\Needspace{18\baselineskip}
Attach two \(206\)-element supports to the two preserved faces. Let

\begin{equation*}
\begin{adjustbox}{max width=\linewidth,center}
$\displaystyle
C=C_{\mathrm{root}}=\{r_0,r_1,r_2\},
\qquad
|C|=3,
$
\end{adjustbox}
\end{equation*}
% LAYOUT_ONLY_GUARD:ROOT_SUPPORT_DEFINITION
\Needspace{13\baselineskip}
be the common \(C_3\) root carrier. Let \(U_{\mathrm{dir}}\) and \(U_{\mathrm{route}}\) be disjoint \(203\)-element sets, each disjoint from \(C\), and define

\begin{equation*}
\begin{adjustbox}{max width=\linewidth,center}
$\displaystyle
\Gamma_{\mathrm{dir}}
=
C\sqcup U_{\mathrm{dir}},
\qquad
\Gamma_{\mathrm{route}}
=
C\sqcup U_{\mathrm{route}}.
$
\end{adjustbox}
\end{equation*}
Then

\begin{equation*}
\begin{adjustbox}{max width=\linewidth,center}
$\displaystyle
|\Gamma_{\mathrm{dir}}|
=
|\Gamma_{\mathrm{route}}|
=
3+203
=206,
\qquad
\Gamma_{\mathrm{dir}}\cap\Gamma_{\mathrm{route}}=C.
$
\end{adjustbox}
\end{equation*}
In plain terms, two 206-element supports are attached to the two preserved faces. They are not \(206\) physical events.

The subtraction term has independent provenance: \(3=|C|\) is the cardinality of the common \(C_3\) root carrier. It is not obtained from \(n_{\mathrm{post}}+n_{\mathrm{route}}=2+1\), despite the numerical equality. The face counts distinguish two aspects of one event; they do not supply the overlap removed in the amalgamation.

\Needspace{5\baselineskip}
\subsection{The Free Pushout}
Form the free pushout

\begin{equation*}
\begin{adjustbox}{max width=\linewidth,center}
$\displaystyle
\Gamma_{\mathrm{scr}}
=
\Gamma_{\mathrm{dir}}
\amalg_C
\Gamma_{\mathrm{route}},
$
\end{adjustbox}
\end{equation*}
which identifies precisely the two copies of each element of \(C\) and introduces no other relation. At least three reductions are required because duplicated common roots would double-count the same-event carrier. At most three are permitted because freeness introduces no relation among the \(406\) non-root elements and non-erasure forbids deleting them. Therefore

\begin{equation}
\label{eq:pushout-409}
\begin{adjustbox}{max width=\linewidth,center}
$\displaystyle
\boxed{
|\Gamma_{\mathrm{scr}}|
=
206+206-|C|
=
206+206-3
=409.
}
$
\end{adjustbox}
\end{equation}
Equation \(\eqref{eq:pushout-409}\) is the complete count; the two face counts do not supply its subtraction term.

The alternative \(412\) double-counts the common roots. The alternative \(406\) requires additional identifications or erasures. Only \(409\) preserves same-event identity, freeness, and non-erasure.

\Needspace{5\baselineskip}
\subsection{The Finite-Screen Factor}
For a completed count \(N\ge3\), equal completed units and one-winding closure force the successor cycle to a regular \(N\)-gon on a unit-diameter circle, up to an overall rotation and carried orientation. Each side has length \(\sin(\pi/N)\), so

\begin{equation*}
\begin{adjustbox}{max width=\linewidth,center}
$\displaystyle
\pi_L(N)=N\sin\frac{\pi}{N},
\qquad
s_N=\frac{N\sin(\pi/N)}{\pi}.
$
\end{adjustbox}
\end{equation*}
At \(N=409\),

\begin{equation}
\label{eq:screen-factor}
\begin{adjustbox}{max width=\linewidth,center}
$\displaystyle
\boxed{
s_{409}
=
\frac{409\sin(\pi/409)}{\pi},
\qquad
0<s_{409}<1.
}
$
\end{adjustbox}
\end{equation}
The exact factor used below is Eq. \(\eqref{eq:screen-factor}\).

\Needspace{5\baselineskip}
\subsection{Observer-Rooted Middle Placement}
The middle-only action is fixed before the finite factor is applied. The lawful ancestry chain and the intrinsic rank chain are

\begin{equation*}
\begin{adjustbox}{max width=\linewidth,center}
$\displaystyle
A\prec B\prec D,
\qquad
\min\prec\operatorname{mid}\prec\max.
$
\end{adjustbox}
\end{equation*}
Here \(D\) is the endpoint label in the ancestry chain; it is not the geometric filled-share invariant denoted by the same letter in Section 4. The geometric use has type

\begin{equation*}
\begin{adjustbox}{max width=\linewidth,center}
$\displaystyle
D_{\mathrm{geom}}=\frac{\pi}{3\sqrt2},
$
\end{adjustbox}
\end{equation*}
whereas the ancestry use is the terminal label in the ordered relational chain. The completed four-cell context and the independently established active coefficient also admit the typed count composite

\begin{equation*}
\begin{adjustbox}{max width=\linewidth,center}
$\displaystyle
4m
=
4\left(\frac1{2\pi}\right)
=
\frac2\pi.
$
\end{adjustbox}
\end{equation*}
This is neither an equality with \(D_{\mathrm{geom}}\) nor a redefinition of either use of \(D\). The two readings meet only at the typed completed four-cell relational location: one describes the geometric filled share, while the other names an ancestry endpoint inside the four-cell observer construction.

The first chain has the unique interior vertex \(B\), and the second has the unique middle rank. A lawful observer readout is a bijective, betweenness-preserving map

\begin{equation*}
\begin{adjustbox}{max width=\linewidth,center}
$\displaystyle
\mathcal R:
\{A,B,D\}
\longrightarrow
\{\min,\operatorname{mid},\max\}.
$
\end{adjustbox}
\end{equation*}
It must therefore carry the unique interior source element to the unique middle target:

\begin{equation}
\label{eq:observer-middle}
\begin{adjustbox}{max width=\linewidth,center}
$\displaystyle
\boxed{
\mathcal R(B)=\operatorname{mid}.
}
$
\end{adjustbox}
\end{equation}
Equation \(\eqref{eq:observer-middle}\) fixes the unique middle placement without changing the intrinsic spectrum.

Reversal may exchange the two endpoints, but it fixes the unique middle. This is observer-rooted asymmetry in its operative form: the intrinsic unordered spectrum is unchanged, while a lawful report depends on the observer's rooted ancestry and betweenness position. The finite-screen action below is therefore selected by observer-rooted betweenness, not by a particle label or by choosing a displayed spectral sector.

The factor acts once, after quadratic-support formation, and only on the middle rank:

\begin{equation*}
\begin{adjustbox}{max width=\linewidth,center}
$\displaystyle
\boxed{
(q_{\min},q_{\mathrm{mid}},q_{\max})
\longmapsto
(q_{\min},s_{409}q_{\mathrm{mid}},q_{\max}).
}
$
\end{adjustbox}
\end{equation*}
% LAYOUT_ONLY_GUARD:STANDALONE_THUS
\Needspace{7\baselineskip}
Thus

\begin{equation*}
\begin{adjustbox}{max width=\linewidth,center}
$\displaystyle
\boxed{
R_{\mathrm{mid}/\min}^{\mathrm{scr}}
=
s_{409}R_{\mathrm{mid}/\min}^{\mathrm{int}},
\qquad
R_{\max/\min}^{\mathrm{scr}}
=
R_{\max/\min}^{\mathrm{int}}.
}
$
\end{adjustbox}
\end{equation*}
Because \(s_{409}>0\), division by it recovers the intrinsic middle ratio. The construction has produced \(206\), \(409\), and \(s_{409}\) without using a measured muon or tau value.

